\subsection{Identification and computational illustrations}
Consider a randomised binary assignment $Z$ to a creative affordance, a proximal mediator $M$ (for example state agency) and a continuous outcome $Y$ oriented so that larger values are favourable. A linear illustrative structural model is
\begin{align}
M &= aZ+U,\\
Y &= \tau Z+bM+cU+\epsilon,
\end{align}
where $U$ is an unobserved mediator--outcome common cause and both disturbances are independent of assignment. For any real $\delta$, define
\[
 b'=b+\delta,\quad c'=c-\delta,\quad \tau'=\tau-a\delta .
\]
Substitution yields the same $M$ and $Y$ for every realisation of $(Z,U,\epsilon)$. Consequently, even randomisation and perfect observation of the mediator do not distinguish these structural models. The total effect $\tau+ab$ is invariant, while the path-specific effect $ab$ changes. This counterexample uses an additive linear model with no treatment--mediator interaction; it demonstrates non-identification in this class, not that every mediation design fails. Randomising a mediator-relevant intervention can test a component-level effect; attribution to the mediator still requires assumptions about other pathways, intervention versions and mediator--outcome confounding.

\paragraph{Goal--affordance contrast.} Let $\mu_{ga}$ denote mean goal attainment under random assignment to goal instruction $g$ and interface $a$, both coded as zero (articulation or fidelity preservation) and one (reinterpretation or perspective transformation). The prespecified equal-cell-weight matching contrast is
\[
 \psi=\frac{\mu_{00}+\mu_{11}-\mu_{01}-\mu_{10}}{2}.
\]
It equals half the conventional difference-in-differences interaction and is positive when matched conditions have higher average attainment than mismatched conditions. Randomise the two factors independently, retain all assigned participants in the estimand, and assess measurement comparability across goal instructions. A positive interaction does not imply that each matched cell individually exceeds its corresponding mismatched cell; report all four means and both simple effects. Confidence intervals and sample size planning must target $\psi$, not the main effect of the interface. Under independent equal-size cells with common residual variance $\sigma^2$ and $n$ participants per cell, $\operatorname{Var}(\hat\psi)=\sigma^2/n$. This planning approximation excludes attrition, clustering and heterogeneous residual variances.

\paragraph{Immediate and delayed effects.} A minimal illustrative response model is
\[
 L=\rho R-\lambda Q,\qquad 0\leq\rho\leq1,\quad \lambda Q\geq0,
\]
where $R$ denotes immediate relief, $L$ delayed relief, $\rho$ persistence and $\lambda Q$ an incremental delayed cost of perseverative thought. For positive immediate relief, delayed deterioration occurs when $\lambda Q>\rho R$. This algebraic example is a possible response trajectory, not a fitted model; other causes of delayed affective change are omitted. The empirical diary model in the validation programme allows additional predictors and participant-level variation.

\paragraph{Reproducible synthetic examples.} A companion Python script illustrates these three properties using arbitrary standardised parameters and a fixed pseudo-random seed. In the first example, models with indirect effects of 0.30 and 0.00 both imply a total effect of 0.40 and produce identical observations up to floating-point rounding. In the factorial example, 1000 simulated observations per cell with $\psi=0.40$ produce $\hat\psi\approx0.382$. Finally, immediate relief of 0.60 and persistence of 0.50 imply delayed relief of 0.30 without perseverative cost and $-0.50$ when that cost is 0.80. These scenarios illustrate structural implications and estimator behaviour; parameter values are not estimates from the target population and should not be used as evidence for efficacy or as an empirical basis for sample size.
